\section{Conclusion}
\label{sec:conclusion}
We introduced ScrapeBench, an evaluation protocol that treats the web acquisition funnel---transcoding, browser and thumbnail rendering, agent capture, restoration, and diffusion-based regeneration---as the primary adversary of face anti-customization defenses. Measuring four published protections, a JPEG-aware retrain, and two controls across three attack families under the funnel's tiers, we found that the headline protection gap is largely an artifact of unmatched baselines---identity-matched controls shrink it roughly three-fold and erase its significance---that what remains erodes slowly under platform transcoding but is attack-conditional, essentially absent for zero-shot extractors that bypass the poisoned fine-tuning gradient, and that the funnel's most expensive tier manufactures pseudo-protection by destroying images. Because funnel transformations cost orders of magnitude less than protection generation and are selected by parties with no stake in the attacker's success, an arms race anchored on stronger perturbations is structurally unwinnable. ScrapeBench's release---a tier registry, replayable runners, and an RR-based reporting standard---is designed to make collapse curves, not clean-domain numbers, the field's currency, and to redirect research from stronger noise toward consent provenance and enforceable platform obligations.
